\documentclass[11pt, letterpaper]{article}
% ----- PAQUETES -----
\usepackage[utf8]{inputenc}
\usepackage[spanish]{babel}
\usepackage{graphicx}
\usepackage[margin=2.3cm]{geometry}
\usepackage{xcolor}
\usepackage{titlesec}
\usepackage{booktabs}
\usepackage{enumitem}
\usepackage{xurl}
\usepackage{hyperref}
\usepackage{fancyhdr}
\usepackage{amsmath}
\usepackage{amssymb}
\usepackage{tikz}
\usetikzlibrary{arrows.meta, positioning, babel}
% ----- COLORES -----
\definecolor{UnalGreen}{RGB}{0,153,51}
\definecolor{DarkGray}{RGB}{60,60,60}
% ----- SECCIONES -----
\titleformat{\section}
{\normalfont\large\bfseries\color{UnalGreen}}
{\thesection}{1em}{}[{\titlerule[0.8pt]}]
\titleformat{\subsection}
{\normalfont\normalsize\bfseries\color{DarkGray}}
{\thesubsection}{1em}{}
\titlespacing*{\section}{0pt}{8pt}{4pt}
\titlespacing*{\subsection}{0pt}{6pt}{2pt}
% ----- ENLACES -----
\hypersetup{
	colorlinks=true,
	linkcolor=UnalGreen,
	urlcolor=UnalGreen
}
% ----- ENCABEZADO -----
\pagestyle{fancy}
\fancyhf{}
\lhead{\footnotesize Fundamentos de Mecánica (1000019) -- Taller}
\rhead{\footnotesize Semestre 2026-2}
\cfoot{\thepage}
% ----- LISTAS COMPACTAS -----
\setlist[itemize]{
	leftmargin=*,
	itemsep=2pt,
	topsep=2pt,
	parsep=0pt
}
\begin{document}
	
	% =========================================================
	% ENCABEZADO
	% =========================================================
	
	\begin{center}
		\vspace*{-0.8cm}
		\includegraphics[width=0.15\textwidth]
		{imagenes/universidad-nacional-de-colombia-sede-bogota-seeklogo.png}\\[0.1cm]
		
		{\Large\bfseries Taller de Cinemática y Graficación -- Grupo 14}\\[0.05cm]
		{\large\bfseries Fundamentos de Mecánica (1000019)}
	\end{center}
	
	
	% =========================================================
	% INFORMACIÓN DEL CURSO
	% ========================================================
	
	\vspace{0.2cm}
	
	% =========================================================
	\section{Cinemática 1D}
	% =========================================================
	
	\begin{enumerate}[label=\arabic*., leftmargin=*, itemsep=8pt]
		
		\item ¿Cuál es la rapidez promedio de un caminante joven, un ciclista aficionado, un auto en la ciudad de Bogotá y un auto de Fórmula 1? Exprese estas cantidades en m/s y en km/h.
		
		\item Use las ecuaciones $x(t) = \frac{1}{2}a_xt^2 + v_{0x}t + x_0$ y $v_x(t) = v_{0x} + a_xt$ para demostrar que
		\[
		v_x^2 = v_{0x}^2 + 2a_x(x-x_0).
		\]
		
		\item Un auto se mueve a lo largo de un camino recto. Antes de $t=2.5$ s viaja a rapidez constante; a partir de ese instante el conductor acelera. La posición del auto en función del tiempo es
		\[
		x(t) =
		\begin{cases}
			8t & 0 \le t \le 2.5 \text{ s} \\[4pt]
			20 + 8(t-2.5) + \dfrac{1}{2}(t-2.5)^2 & t > 2.5\text{ s}
		\end{cases}
		\qquad [x] = \text{m},\ [t]=\text{s}.
		\]
		\begin{center}
			\begin{tikzpicture}[x=0.85cm,y=0.055cm]
				% ----- CUADRÍCULA (grid) -----
				\draw[gray!20, very thin] (0,-2) grid[xstep=2.5, ystep=20] (9.5,90);
				
				% ----- EJES -----
				\draw[->, thick] (0,0) -- (10.5,0) node[right]{\normalsize $t$ (s)};
				\draw[->, thick] (0,0) -- (0,98) node[above]{\normalsize $x$ (m)};
				
				% ----- TICKS EN X -----
				\foreach \tk in {0, 2.5, 5, 7.5, 10} {
					\draw (\tk,0) -- (\tk,-2.5) node[below, yshift=-2pt]{\normalsize \tk};
				}
				
				% ----- TICKS EN Y (CORREGIDOS: líneas más cortas) -----
				\foreach \yk in {0, 20, 40, 60, 80} {
					\draw (0,\yk) -- (-0.3,\yk) node[left, xshift=-3pt]{\normalsize \yk};
				}
				
				% ----- FUNCIONES -----
				% Línea punteada (sin acelerar)
				\draw[thick, dashed, DarkGray] plot[domain=2.5:9.3, samples=2] (\x,{8*\x});
				
				% Tramo 1: velocidad constante
				\draw[thick, UnalGreen] plot[domain=0:2.5, samples=15] (\x,{8*\x});
				
				% Tramo 2: acelerado
				\draw[thick, UnalGreen] plot[domain=2.5:9.3, samples=60] (\x,{20+8*(\x-2.5)+0.5*(\x-2.5)*(\x-2.5)});
				
				% ----- PUNTO CLAVE t=2.5 s -----
				\fill (2.5,20) circle (2pt);
				
			\end{tikzpicture}\\
			\small\textit{Figura 1: posición del auto en función del tiempo (la línea punteada muestra la trayectoria que habría seguido sin acelerar).}
		\end{center}
		\begin{enumerate}[label=\alph*)]
			\item ¿Cuál es la velocidad del auto antes de $t=2.5$ s?
			\item Calcule la velocidad media entre $t_1=6$ s y $t_2=7.5$ s.
			\item Calcule la velocidad media entre $t_1=7$ s y $t_2=7.5$ s.
			\item Calcule la velocidad media entre $t_1=7.45$ s y $t_2=7.5$ s.
			\item ¿Qué observa que ocurre con estos valores a medida que el intervalo se hace más pequeño? ¿Cuál es la velocidad instantánea en $t=7.5$ s?
		\end{enumerate}
		
		\item Un jugador de béisbol, después de ganar un partido, se quita la gorra y la lanza verticalmente hacia arriba con una rapidez inicial de $10$ m/s. La gorra experimenta una desaceleración de $9.8$ m/s$^2$.
		\begin{enumerate}[label=\alph*)]
			\item ¿Cuánto tiempo le toma a la gorra alcanzar su altura máxima?
			\item ¿Cuál es esa altura máxima?
			\item Si la gorra es atrapada a la misma altura desde la que fue lanzada, ¿cuál es su tiempo total de vuelo?
		\end{enumerate}
		
		\item Un auto se mueve a lo largo del eje $x$ con velocidad $v_x$ dada por la Figura 2. Si $x_0=2$ m en $t_0=2$ s, determine la posición del auto en $t=10$ s.
		\begin{center}
			\begin{tikzpicture}[x=0.62cm,y=0.17cm]
				% ----- CUADRÍCULA (grid) -----
				\foreach \tk in {0,2,4,6,8,10} \draw[gray!20, very thin] (\tk,-8) -- (\tk,24);
				\foreach \vk in {-8,-4,0,4,8,12,16,20,24} \draw[gray!20, very thin] (0,\vk) -- (10,\vk);
				
				% ----- EJES -----
				\draw[->, thick] (0,-8) -- (0,26) node[above]{\normalsize $v_x$ (m/s)};
				\draw[->, thick] (-0.3,0) -- (10.4,0) node[right]{\normalsize $t$ (s)};
				
				% ----- TICKS EN X (más cortos) -----
				\foreach \tk in {2,4,6,8,10} {
					\draw[thick] (\tk,0) -- (\tk,-0.3) node[below, yshift=-4pt]{\normalsize \tk};
				}
				
				% ----- TICKS EN Y (más cortos) -----
				\foreach \vk in {-8,-4,0,4,8,12,16,20,24} {
					\draw[thick] (0,\vk) -- (-0.3,\vk) node[left, xshift=-4pt]{\normalsize \vk};
				}
				
				% ----- FUNCIÓN -----
				\draw[thick, UnalGreen] (0,20) -- (6,-4) -- (10,4);
				
			\end{tikzpicture}\\
			\small\textit{Figura 2: velocidad del auto en función del tiempo.}
		\end{center}
		\textit{Sugerencia: recuerden que, como se vio en la Sesión 2, el área bajo la curva $v_x$ vs.\ $t$ corresponde al desplazamiento (suma de Riemann).}
		
		\item La posición de una partícula que se mueve en línea recta está dada por $x(t) = t^3 - 9t^2 + 15t + 2$ (m, con $t$ en s).
		\begin{enumerate}[label=\alph*)]
			\item Encuentre $v_x(t)$ y determine analíticamente los instantes en que la rapidez de la partícula es cero.
			\item Verifiquen su resultado identificando esos mismos puntos en la gráfica de $x(t)$ (Figura 3).
			\item ¿La partícula cambia de sentido de movimiento en alguno de esos instantes? Justifiquen.
		\end{enumerate}
		\begin{center}
			\begin{tikzpicture}[x=1.0cm,y=0.13cm]
				% ----- CUADRÍCULA (grid) -----
				\draw[gray!20, very thin] (0,-26) grid[xstep=1, ystep=5] (7,12);
				
				% ----- EJES -----
				\draw[->, thick] (0,-26) -- (0,13) node[above]{\normalsize $x$ (m)};
				\draw[->, thick] (-0.3,0) -- (7.4,0) node[right]{\normalsize $t$ (s)};
				
				% ----- TICKS EN X (más cortos) -----
				\foreach \tk in {1,2,3,4,5,6,7} {
					\draw[thick] (\tk,0) -- (\tk,-0.7) node[below, yshift=-4pt]{\normalsize \tk};
				}
				
				% ----- TICKS EN Y (más cortos) -----
				\foreach \yk in {-25,-20,-15,-10,-5,0,5,10} {
					\draw[thick] (0,\yk) -- (-0.1,\yk) node[left, xshift=-4pt]{\normalsize \yk};
				}
				
				% ----- FUNCIÓN -----
				\draw[thick, UnalGreen] plot[domain=0:7, samples=80] (\x,{\x*\x*\x-9*\x*\x+15*\x+2});
				
				% ----- PUNTOS DESTACADOS -----
				\fill[red!70!black] (1,9) circle (2.5pt) node[above right, yshift=2pt, xshift=2pt]{\normalsize $(1,9)$};
				\fill[red!70!black] (5,-23) circle (2.5pt) node[below right, yshift=-2pt, xshift=2pt]{\normalsize $(5,-23)$};
				
			\end{tikzpicture}\\
			\small\textit{Figura 3: posición de la partícula en función del tiempo.}
		\end{center}
		
	\end{enumerate}
	
	\section{Graficación (OPCIONAL)}
	
	En esta actividad van a usar un simulador virtual para medir tiempos de caída, organizar sus datos y graficarlos a mano. El objetivo es obtener experimentalmente el valor de $g$ (gravedad) usando la misma idea que Galileo: graficar altura vs. $t^2$ para obtener una línea recta.
	
	\subsection{Paso 1: Tomar los datos}
	
	\begin{enumerate}
		\item Abran el simulador de caída libre:
		
		\url{https://aulaquest.github.io/simulador-caida-libre/}
		
		(Si conocen otro, pueden usarlo, pero este permite exportar datos en \textbf{CSV}, que significa \emph{valores separados por comas}).
		
		\item Configúrenlo así:
		\begin{itemize}
			\item Sin rozamiento.
			\item Velocidad inicial cero.
			\item Planeta Tierra.
		\end{itemize}
		
		\item Suelten el objeto desde \textbf{al menos 5 alturas distintas} (por ejemplo: $2, 4, 6, 8, 10$ m).
		
		\item Para cada altura, anoten el tiempo de caída. Si usan el simulador recomendado, pueden descargar el archivo CSV; el tiempo de caída es el \textbf{último valor de la columna de tiempo} (cuando la posición es cero).
		
		\item Llenen esta tabla con sus datos:
	\end{enumerate}
	
	\begin{center}
		\begin{tabular}{|c|c|c|}
			\hline
			Altura $h$ (m) & Tiempo $t$ (s) & $t^2$ (s$^2$) \\
			\hline
			& & \\
			& & \\
			& & \\
			& & \\
			& & \\
			\hline
		\end{tabular}
	\end{center}
	
	\subsection{Paso 2: Graficar a mano}
	
	\begin{enumerate}
		\item En papel milimetrado (o en una hoja cuadriculada), dibujen un plano cartesiano:
		\begin{itemize}
			\item Eje horizontal ($x$): $t^2$ (s$^2$).
			\item Eje vertical ($y$): altura $h$ (m).
		\end{itemize}
		
		\item Ubiquen sus puntos $(t^2, h)$.
		
		\item Dibujen \textbf{la mejor línea recta} que pase lo más cerca posible de todos los puntos. No tiene que pasar exactamente por todos, solo debe representar la tendencia general.
	\end{enumerate}
	
	\subsection{Paso 3: Calcular $g$}
	
	La física dice que:
	\[
	h = \frac{1}{2} g \, t^2
	\]
	Eso significa que la pendiente de su recta es $m = g/2$. Por lo tanto:
	\[
	g_{\text{exp}} = 2 \times \text{pendiente}
	\]
	
	Para calcular la pendiente, elijan dos puntos sobre la recta (no necesariamente de la tabla) y usen:
	\[
	m = \frac{\Delta h}{\Delta t^2}
	\]
	
	\subsection{Paso 4: Análisis rápido}
	
	\begin{enumerate}
		\item Comparen su $g_{\text{exp}}$ con $9.8$ m/s$^2$ usando:
		\[
		\text{Error \%} = \frac{|g_{\text{exp}} - 9.8|}{9.8} \times 100
		\]
		
		\item ¿Los puntos quedaron más o menos alineados? ¿Por qué creen que sí o no?
		
		\item Mencionen dos posibles fuentes de error (ej: leer el tiempo en la pantalla, redondeo, etc.).
	\end{enumerate}
	
	\vspace{1cm}
	\noindent
	\textbf{Importante:} No se preocupen por la perfección. Lo importante es practicar el proceso de graficar, linealizar y extraer información de una gráfica.
	
	\begin{center}
		\small\textit{*Ejercicios tomados y adaptados del taller de la profesora Edna Carolina Pinilla Beltrán.}
	\end{center}
	
	\vspace{0.5cm}
	{\color{UnalGreen}\hrule height 0.5pt}
	\begin{center}
		\small 
		\textit{Trabajo en grupos.}
	\end{center}
	
\end{document}